(1) 一个热电制冷元件冷端温度 \(T_c>0\)、热端温度 \(T_h>T_c\)，Seebeck系数差 \(\alpha>0\)、总内阻 \(R>0\)、被动热导 \(K>0\)。电流 \(I>0\) 方向定义为冷端吸热。假定焦耳热各一半回到两端，列出冷端吸热率 \(\dot Q_c\)。
(2) 求热端放热率 \(\dot Q_h\) 及供电功率 \(P_{\rm el}\)，验证第一定律。
(3) 在温差固定时求使 \(\dot Q_c\) 最大的电流及最大制冷量。
(4) 给出制冷可能的阈值条件与工作制冷系数 COP，并讨论忽略焦耳热导致的错误。
